\documentclass[10pt,a4paper]{amsart}
\usepackage{amsmath,amssymb,amsthm,graphicx}
% Retained minimal preamble for the Stacks Project chapter more-algebra.tex (More on Algebra).
% The chapter file pinned-original-source.tex carries only \input{preamble} on its first line, so no
% preamble text lives inside the chapter itself.  The Stacks Project uses one single shared preamble for
% every chapter; the snapshot stages it identically as sources/stacks-chow/preamble.tex and
% sources/stacks-etale-cohomology/preamble.tex.  The block below is copied mechanically, line for line,
% from sources/stacks-chow/preamble.tex lines 14-18 (the Xy-pic setup) and 154-217 (the native theorem
% declarations together with the mathematical macros).  Nothing was added, removed or redefined; the
% wrapper class and the amsmath/amssymb/amsthm/graphicx base packages are supplied by the assembler.
% Chapter: more-algebra.tex -- More on Algebra.
% For commutative diagrams we use Xy-pic
\usepackage[all]{xy}

% We use 2cell for 2-commutative diagrams.
\xyoption{2cell}

% Theorem environments.
%
\theoremstyle{plain}
\newtheorem{theorem}[subsection]{Theorem}
\newtheorem{proposition}[subsection]{Proposition}
\newtheorem{lemma}[subsection]{Lemma}

\theoremstyle{definition}
\newtheorem{definition}[subsection]{Definition}
\newtheorem{example}[subsection]{Example}
\newtheorem{exercise}[subsection]{Exercise}
\newtheorem{situation}[subsection]{Situation}

\theoremstyle{remark}
\newtheorem{remark}[subsection]{Remark}
\newtheorem{remarks}[subsection]{Remarks}

\numberwithin{equation}{subsection}

% Macros
%
\def\lim{\mathop{\mathrm{lim}}\nolimits}
\def\colim{\mathop{\mathrm{colim}}\nolimits}
\def\Spec{\mathop{\mathrm{Spec}}}
\def\Hom{\mathop{\mathrm{Hom}}\nolimits}
\def\Ext{\mathop{\mathrm{Ext}}\nolimits}
\def\SheafHom{\mathop{\mathcal{H}\!\mathit{om}}\nolimits}
\def\SheafExt{\mathop{\mathcal{E}\!\mathit{xt}}\nolimits}
\def\Sch{\mathit{Sch}}
\def\Mor{\mathop{\mathrm{Mor}}\nolimits}
\def\Ob{\mathop{\mathrm{Ob}}\nolimits}
\def\Sh{\mathop{\mathit{Sh}}\nolimits}
\def\NL{\mathop{N\!L}\nolimits}
\def\CH{\mathop{\mathrm{CH}}\nolimits}
\def\proetale{{pro\text{-}\acute{e}tale}}
\def\etale{{\acute{e}tale}}
\def\QCoh{\mathit{QCoh}}
\def\Ker{\mathop{\mathrm{Ker}}}
\def\Im{\mathop{\mathrm{Im}}}
\def\Coker{\mathop{\mathrm{Coker}}}
\def\Coim{\mathop{\mathrm{Coim}}}

% Boxtimes
%
\DeclareMathSymbol{\boxtimes}{\mathbin}{AMSa}{"02}

%
% Macros for moduli stacks/spaces
%
\def\QCohstack{\mathcal{QC}\!\mathit{oh}}
\def\Cohstack{\mathcal{C}\!\mathit{oh}}
\def\Spacesstack{\mathcal{S}\!\mathit{paces}}
\def\Quotfunctor{\mathrm{Quot}}
\def\Hilbfunctor{\mathrm{Hilb}}
\def\Curvesstack{\mathcal{C}\!\mathit{urves}}
\def\Polarizedstack{\mathcal{P}\!\mathit{olarized}}
\def\Complexesstack{\mathcal{C}\!\mathit{omplexes}}
% \Pic is the operator that assigns to X its picard group, usage \Pic(X)
% \Picardstack_{X/B} denotes the Picard stack of X over B
% \Picardfunctor_{X/B} denotes the Picard functor of X over B
\def\Pic{\mathop{\mathrm{Pic}}\nolimits}
\def\Picardstack{\mathcal{P}\!\mathit{ic}}
\def\Picardfunctor{\mathrm{Pic}}
\def\Deformationcategory{\mathcal{D}\!\mathit{ef}}

% --- appended 2026-10-12 (Chief): the project own definitions of the reference/slogan/history
% display environments, copied mechanically from the shared Stacks preamble lines 8-12
% (\usepackage{verbatim} plus the three \newenvironment lines, which the chapter preamble keeps
% verbatim from the project preamble). Nothing added, removed or redefined: these environments
% carry no mathematical content and the project itself typesets them as comments. Without these
% three lines a counted statement whose own text uses them fails with Environment slogan undefined.
\usepackage{verbatim}
\newenvironment{reference}{\comment}{\endcomment}
\newenvironment{slogan}{\comment}{\endcomment}
\newenvironment{history}{\comment}{\endcomment}
\usepackage{cleveref}
\crefname{originallemma}{original lemma}{original lemmas}
\Crefname{originallemma}{Original lemma}{Original lemmas}
\makeatletter
\newlabel{more-algebra-theorem-fs-regular}{{0H7U}{source}{}{sourcepointer-1}{}}
\newlabel{more-algebra-theorem-fs-regular@cref}{{[originallemma][0][]0H7U}{[1][0][]source}{}{}{}}
\makeatother
\crefname{originallemma}{original lemma}{original lemmas}
\Crefname{originallemma}{Original lemma}{Original lemmas}
\makeatletter
\newlabel{section-G-ring}{{07GG}{source}{}{sourcepointer-2}{}}
\newlabel{section-G-ring@cref}{{[originallemma][0][]07GG}{[1][0][]source}{}{}{}}
\makeatother
\crefname{originallemma}{original lemma}{original lemmas}
\Crefname{originallemma}{Original lemma}{Original lemmas}
\makeatletter
\newlabel{section-excellent}{{07QS}{source}{}{sourcepointer-3}{}}
\newlabel{section-excellent@cref}{{[originallemma][0][]07QS}{[1][0][]source}{}{}{}}
\makeatother
\crefname{originallemma}{original lemma}{original lemmas}
\Crefname{originallemma}{Original lemma}{Original lemmas}
\makeatletter
\newlabel{proposition-fs-flat-fibre-fs}{{07NQ}{source}{}{sourcepointer-4}{}}
\newlabel{proposition-fs-flat-fibre-fs@cref}{{[originallemma][0][]07NQ}{[1][0][]source}{}{}{}}
\makeatother
\crefname{originallemma}{original lemma}{original lemmas}
\Crefname{originallemma}{Original lemma}{Original lemmas}
\makeatletter
\newlabel{lemma-formally-smooth-completion}{{07ED}{source}{}{sourcepointer-5}{}}
\newlabel{lemma-formally-smooth-completion@cref}{{[originallemma][0][]07ED}{[1][0][]source}{}{}{}}
\makeatother
\crefname{originallemma}{original lemma}{original lemmas}
\Crefname{originallemma}{Original lemma}{Original lemmas}
\makeatletter
\newlabel{proposition-fs-regular}{{07PM}{source}{}{sourcepointer-6}{}}
\newlabel{proposition-fs-regular@cref}{{[originallemma][0][]07PM}{[1][0][]source}{}{}{}}
\makeatother
\crefname{originallemma}{original lemma}{original lemmas}
\Crefname{originallemma}{Original lemma}{Original lemmas}
\makeatletter
\newlabel{lemma-regular-composition}{{07QI}{source}{}{sourcepointer-7}{}}
\newlabel{lemma-regular-composition@cref}{{[originallemma][0][]07QI}{[1][0][]source}{}{}{}}
\makeatother
\crefname{originallemma}{original lemma}{original lemmas}
\Crefname{originallemma}{Original lemma}{Original lemmas}
\makeatletter
\newlabel{lemma-regular-permanence}{{07NT}{source}{}{sourcepointer-8}{}}
\newlabel{lemma-regular-permanence@cref}{{[originallemma][0][]07NT}{[1][0][]source}{}{}{}}
\makeatother
\crefname{originalreference}{original reference}{original references}
\Crefname{originalreference}{Original reference}{Original references}
\makeatletter
\newlabel{Andre-smooth}{{cite-Andre-smooth}{source}{}{sourcepointer-9}{}}
\newlabel{Andre-smooth@cref}{{[originalreference][0][]cite-Andre-smooth}{[1][0][]source}{}{}{}}
\makeatother
\title{Andre's theorem: regularity, flatness with geometrically regular fibre, and formal smoothness of a local map}
\author{The Stacks Project Authors}
\date{stacks-project:more-algebra.tex:local-snapshot-20261012; GNU FDL 1.2 or later}
\begin{document}
\maketitle
\setcounter{subsection}{1}
\noindent\textit{Editorial note.} All delivered mathematical text is the author's own, copied from the pinned Stacks Project chapter file more-algebra.tex; no statement or proof step was rewritten, reordered or removed.\par\medskip
\noindent\textit{Editorial note.} The chapter file pinned-original-source.tex carries only \textbackslash{}input\{preamble\} on its first line, so no preamble text lives inside the chapter. The retained file sources/stacks-more-algebra/retained-preamble.tex is a mechanical, line-for-line copy of the project's single shared preamble as staged in the snapshot (Xy-pic setup plus theorem declarations and mathematical macros). Nothing was added, removed or redefined, and the excerpt serves only as preamble, never as a span coordinate.\par\medskip
\noindent\textit{Editorial note.} Both counted spans start at the \textbackslash{}begin line of their environment, so the statement and the author's proof are delivered whole; surrounding prose of the section is deliberately outside the counted spans and is carried as standard prerequisites instead.\par\medskip
\noindent\textit{Editorial note.} The citation fields (source\_id, source\_version, authors, source\_url, license.name, source\_title) are identical to those of every other card this lane has produced from this chapter, so the batch selector treats them as one source.\par\medskip
\noindent\textit{Editorial note.} Benign-screen note: the proof contains the phrase 'it is clear that' / 'is obvious' once, used as a legitimate step inside a complete argument; the author proves the statement, so this is not a hint-only proof.\par\medskip
\noindent\textit{Allowed standard prerequisites (editorial).}\begin{itemize}
\item Equivalence of flatness with geometrically regular fibre, flatness with formally smooth fibre, and formal smoothness of the map. Proposition \textbackslash{}ref\{proposition-fs-flat-fibre-fs\}, cited at proof lines 12141-12142. Gives the equivalence of the last three conditions.
\item Regularity implies flatness with geometrically regular fibre. used at proof lines 12143-12144. Gives one of the implications.
\item Andre's theorem on formal smoothness implying regularity. the reference environment citing \textbackslash{}cite\{Andre-smooth\}, cited at proof lines 12119-12121. Supplies the substantive implication; the reference environment is a comment environment in this chapter.
\end{itemize}\medskip
\ifdefined\hypertarget\hypertarget{sourcepointer-1}{}\fi
\noindent\textit{Source reference.} Original lemma 0H7U: focused-seeds/AGP/sources/stacks-snapshot-20261011/tags.txt, source lines 21469--21469. The lemma's own canonical tag, quoted so that readers can resolve the label inside the delivered text; it points at the counted statement itself and carries no proof content.\par\medskip
\ifdefined\hypertarget\hypertarget{sourcepointer-2}{}\fi
\noindent\textit{Source reference.} Original lemma 07GG: focused-seeds/AGP/sources/stacks-snapshot-20261011/tags.txt, source lines 9516--9516. Navigation pointer to a statement the counted proof refers to but does not contain; nothing load-bearing is delegated to it.\par\medskip
\ifdefined\hypertarget\hypertarget{sourcepointer-3}{}\fi
\noindent\textit{Source reference.} Original lemma 07QS: focused-seeds/AGP/sources/stacks-snapshot-20261011/tags.txt, source lines 9842--9842. Navigation pointer to a statement the counted proof refers to but does not contain; nothing load-bearing is delegated to it.\par\medskip
\ifdefined\hypertarget\hypertarget{sourcepointer-4}{}\fi
\noindent\textit{Source reference.} Original lemma 07NQ: focused-seeds/AGP/sources/stacks-snapshot-20261011/tags.txt, source lines 9770--9770. Navigation pointer to a statement the counted proof refers to but does not contain; nothing load-bearing is delegated to it.\par\medskip
\ifdefined\hypertarget\hypertarget{sourcepointer-5}{}\fi
\noindent\textit{Source reference.} Original lemma 07ED: focused-seeds/AGP/sources/stacks-snapshot-20261011/tags.txt, source lines 9443--9443. Navigation pointer to a statement the counted proof refers to but does not contain; nothing load-bearing is delegated to it.\par\medskip
\ifdefined\hypertarget\hypertarget{sourcepointer-6}{}\fi
\noindent\textit{Source reference.} Original lemma 07PM: focused-seeds/AGP/sources/stacks-snapshot-20261011/tags.txt, source lines 9802--9802. Navigation pointer to a statement the counted proof refers to but does not contain; nothing load-bearing is delegated to it.\par\medskip
\ifdefined\hypertarget\hypertarget{sourcepointer-7}{}\fi
\noindent\textit{Source reference.} Original lemma 07QI: focused-seeds/AGP/sources/stacks-snapshot-20261011/tags.txt, source lines 9833--9833. Navigation pointer to a statement the counted proof refers to but does not contain; nothing load-bearing is delegated to it.\par\medskip
\ifdefined\hypertarget\hypertarget{sourcepointer-8}{}\fi
\noindent\textit{Source reference.} Original lemma 07NT: focused-seeds/AGP/sources/stacks-snapshot-20261011/tags.txt, source lines 9773--9773. Navigation pointer to a statement the counted proof refers to but does not contain; nothing load-bearing is delegated to it.\par\medskip
\ifdefined\hypertarget\hypertarget{sourcepointer-9}{}\fi
\noindent\textit{Source reference.} Original reference cite-Andre-smooth: sources/stacks-more-algebra/pinned-original-source.tex, source lines 12118--12118. Bibliographic pointer appearing in the counted text; it carries no mathematical argument.\par\medskip
\begin{theorem}[Andr\'e]
\label{theorem-fs-regular}
\begin{reference}
The implication (4) $\Rightarrow$ (1) is the main result of \cite{Andre-smooth}.
\end{reference}
Let $A \to B$ be a local homomorphism of Noetherian local rings.
Let $k$ be the residue field of $A$ and $\overline{B} = B \otimes_A k$
the special fibre. Assume $A \to A^\wedge$ is regular\footnote{For example,
if $A$ is a G-ring or (quasi-)excellent, see Sections \ref{section-G-ring} and
\ref{section-excellent}.}.
The following are equivalent
\begin{enumerate}
\item $A \to B$ is regular,
\item $A \to B$ is flat and $\overline{B}$ is geometrically regular
over $k$,
\item $A \to B$ is flat and $k \to \overline{B}$ is formally smooth
in the $\mathfrak m_{\overline{B}}$-adic topology, and
\item $A \to B$ is formally smooth in the $\mathfrak m_B$-adic
topology.
\end{enumerate}
\end{theorem}

\begin{proof}
We have seen the equivalence of (2), (3), and (4) in
Proposition \ref{proposition-fs-flat-fibre-fs}.
It is clear that (1) implies (2). Thus we assume (4) holds
and we prove (1).

\medskip\noindent
By Lemma \ref{lemma-formally-smooth-completion} we see that
$A^\wedge \to B^\wedge$ is formally smooth. By
Proposition \ref{proposition-fs-regular} we see that $A^\wedge \to B^\wedge$
is regular. By assumption $A \to A^\wedge$ is regular, hence $A \to B^\wedge$
is regular by Lemma \ref{lemma-regular-composition}.
Since $B \to B^\wedge$ is faithfully flat, we conclude that
$A \to B$ is regular by Lemma \ref{lemma-regular-permanence}.
\end{proof}

\end{document}
